\section{Discussion}

\subsection{Finding stable spinodal compositions based on the Hessian of the demixing free energy}

The stability of spinodal demixing within the studied body-centered cubic iron-chrome alloy system was investigated through the second derivative of the free energy with respect to chromium composition. The analysis revealed that spinodal demixing is stable between 0.26$\pm$0.01 and 0.95$\pm$0.02 chromium content \cite{korzhavyi2009electronic,trochet2021influence,levesque2011simple, binder2001spinodal}. This finding is significant as it defines the compositional boundaries within which spinodal decomposition is thermodynamically favored.

\subsubsection{Minimizing elastic anisotropy through spinodal demixing}
Elastic anisotropy in polycrystalline alloys can lead to preferential mechanical behavior along certain crystallographic directions, which is often undesirable for isotropic mechanical properties. Based on the work of Kádas et al. \cite{kadas2011thermo} and the findings from the European Research Council's Final Report Summary \cite{euroresearchcouncil2014}, the reduction of elastic anisotropy through spinodal demixing has been explored. Within the spinodal bounds, the demixing process appears to be a mechanism through which the material minimizes its free energy, consequently reducing the elastic anisotropy in the alloy system.

The impact of chromium content on the polycrystalline anisotropy of body-centered cubic iron-chrome alloys was quantified. It was observed that increasing chromium content within the spinodal bounds leads to a decrease in elastic anisotropy, which in turn dynamically stabilizes the alloy. This behavior underscores the importance of compositional control in tailoring the mechanical properties of alloys.


\subsection{Detecting thermally-induced spontaneous demixing towards embrittlement}
Both Pugh's classical criterion and Pettifor's quantum criterion have been discussed in the context of the ductile-brittle transition. These approaches provide complementary perspectives on the mechanical behavior of materials. The method proposed for investigating demixing and embrittlement phenomena is not only effective but also scalable to high-throughput datasets, as indicated by Senkov and Miracle \cite{senkov2021}. The application of this method to large-scale materials datasets holds the potential to expedite the discovery of new materials with tailored properties.

Finally, the inclusion of validation measurements is critical to ensure the reliability and accuracy of the results. Such measurements provide confidence in the predictive capabilities of the computational models and the reproducibility of the findings. This multi-faceted approach, combining theoretical predictions with experimental validation, paves the way for a more comprehensive understanding of the complex phenomena governing material stability and mechanical behavior.